% SHARD-ID: CA-29-GAUSSIAN-BOSON-REAL-SPACE-DENSITY
% SHARD-TITLE: Gaussian Boson Real-Space Energy Density
% SHARD-SUMMARY: Fixes the scalar Gaussian real-space energy-density convention as a local split of the sourced lattice Hamiltonian.
% SHARD-SUMMARY: Records cell-volume, equal bond-sharing, boundary, vacuum-subtraction, and improvement policies before position-weighted use.
% SHARD-KEYWORDS: Gaussian boson, real-space density, Hamiltonian density, cell volume, bond sharing, normal ordering, improvement

\section{Gaussian Boson Real-Space Energy Density}
\label{sec:gaussian-boson-real-space-density}

\subsection{Status, Sources, and Dependencies}

\textbf{Status: convention, not a theorem.}  This shard fixes how the scalar
Gaussian Hamiltonian is split into site-attached real-space energy densities.
It preserves the total Hamiltonian used in the sourced nearest-neighbour
scalar model and in CA-24--CA-28, but it does not claim a continuum tensor
object, a conservation law, or covariance.

Dependencies: CA-23--CA-28 and CONVENTIONS.md (j), (l).

% Source: literature/md/2010.11121/2010.11121.md:598--603 -- sourced lattice scalar Hamiltonian with cell-volume factor and dispersion.
% Source: literature/md/2010.11121/2010.11121.md:614--623 -- physical mass convention and renormalized lattice Klein-Gordon dispersion.
% Source: references/text/CFTFromLatticeFermions.txt:378--389 -- vacuum/ground-state subtraction corresponds to normal ordering in the free-fermion Koo-Saleur setting.
% Convention: CONVENTIONS.md (j) -- scalar Gaussian coefficient tier and finite periodic symbol/eigenvalue scope.
% Convention: CONVENTIONS.md (l) -- real-space energy-density split fixed here.

\subsection{Cell Volume and Notation}

Let the hypercubic lattice spacing be \(\epsilon\), so one spatial cell has
volume \(\epsilon^d\).  In this shard \(h_x\) is the physical energy density
attached to the cell at \(x\):
\[
  H_\epsilon=\epsilon^d\sum_x h_x.
\]
The integrated site energy is
\[
  e_x:=\epsilon^d h_x,
  \qquad
  H_\epsilon=\sum_x e_x .
\]
This distinction matters because the generic first-moment notation of
CONVENTIONS.md (h) writes \(H=\sum_x h_x\); for the Gaussian block, the object
to insert there is \(e_x\), not the density \(h_x\).

For fields \(\Phi_x,\Pi_x\) and a real symmetric translation-invariant kernel
\(V_r=V_{-r}\), use
\[
  H_\epsilon
  =
  \epsilon^d
  \left(
    \frac12\sum_x\Pi_x^2
    +
    \frac12\sum_{x,r}\Phi_xV_r\Phi_{x+r}
  \right).
\]
For the nearest-neighbour Klein-Gordon source, the renormalized coefficients are
\[
  V_0=m^2+\frac{2d}{\epsilon^2},
  \qquad
  V_{\pm e_a}=-\frac{1}{\epsilon^2},
\]
which give the dispersion recorded in CA-24.

\subsection{Equal-Endpoint Bond Split}

The chosen local density is
\[
  h_x
  =
  \frac12\Pi_x^2
  +
  \frac12 V_0\Phi_x^2
  +
  \frac14\sum_{r\ne0}V_r
  \bigl(\Phi_x\Phi_{x+r}+\Phi_{x-r}\Phi_x\bigr).
\]
Each nonzero displacement bond is assigned half to each endpoint cell.  Summing
over \(x\) gives
\[
  \sum_x h_x
  =
  \frac12\sum_x\Pi_x^2
  +
  \frac12\sum_{x,r}\Phi_xV_r\Phi_{x+r},
\]
so this convention reproduces the total scalar Hamiltonian above after the
cell-volume factor is restored.

The \(q,p\) notation in CA-24--CA-28 absorbs the cell volume into the coefficient
bookkeeping.  In that notation the same displayed formula defines \(e_x\), and
the physical density is \(h_x=e_x/\epsilon^d\).

\subsection{Boundary Policy}

The default exact identities are infinite-lattice identities with summable
tails, or finite periodic total-energy identities.  On a finite open region,
the density is obtained from the finite Hamiltonian actually named: include
only bonds present in that Hamiltonian, and share each included bond equally
between its included endpoints.  No exterior ghost sites, omitted half-bonds,
or boundary counterterms are implicit.

For periodic finite lattices, wrap-around bonds are included in the total
energy split.  This does not create a periodic position convention.  Any
periodic first moment still needs the branch, sawtooth, or interpolation policy
required by CONVENTIONS.md (j).

\subsection{Vacuum-Energy Policy}

The density above is bare.  No normal-ordering or vacuum subtraction is
implicit.  Once a state \(\omega\) is named, the expectation-subtracted version
may be recorded as
\[
  {:}h_x{:}_\omega:=h_x-\omega(h_x)\mathbf 1 .
\]
This changes \(H_\epsilon\) by a scalar and therefore drops out of commutators,
but it matters for absolute energies, expectation values, and Fourier modes of
quadratic densities.  The local subtraction source is the free-fermion
Koo-Saleur setting, where the scaling-limit representation relies on subtracting
vacuum or ground-state expectation values; this shard only records the analogous
policy needed before such quantities are used in the scalar Gaussian block.

\subsection{Improvement Policy}

The default convention has no divergence improvement:
\[
  h'_x=h_x
\]
unless a later shard explicitly names an edge field \(b\), a discrete divergence
\((\nabla_\epsilon\cdot b)_x\), and a boundary convention.  Such a rewrite
\[
  h'_x=h_x+(\nabla_\epsilon\cdot b)_x
\]
may leave the total energy unchanged on a periodic lattice or for decaying
fields, but it is not harmless for position-weighted generators.  A discrete
summation-by-parts calculation shifts first moments by an edge-field term plus
boundary contributions.  For example, in one dimension with
\[
  (\nabla_\epsilon b)_n
  =
  \frac{b_{n+1/2}-b_{n-1/2}}{\epsilon},
  \qquad x_n=\epsilon n,
\]
one has
\[
  \epsilon\sum_n x_n(\nabla_\epsilon b)_n
  =
  -\epsilon\sum_n b_{n+1/2}
  +\text{boundary terms}.
\]
Therefore improvements are separate data, not gauge choices to be silently
modded out.

\subsection{Compiler Consequence}

A scalar Gaussian real-space generator proposal must now carry the tuple
\[
\bigl(\epsilon^d,\; h_x,\; e_x=\epsilon^dh_x,\;
  \text{boundary policy},\; \text{subtraction policy},\;
  \text{improvement policy}\bigr).
\]
The total symbol \(\omega(k)^2\) remains sufficient for CA-24--CA-28, but it is
not sufficient once a construction uses position weights or local quadratic
modes.  Those later constructions must cite this density convention or perform
an explicit convention sweep.
